\input{preamble}

\title{Section 6.6 --- Function Transformations - Answer Key}

\NewDocumentEnvironment{fpa}{m m m m m >{\SplitArgument{3}{,}}D(){5,5,5,5} O{1} +b}{%
a. $f(x)=#1$ (#2)\\
b. $g(x)=#3$\\
c. \begin{GridFitFilled}[#7]#6(2in)
#8
\end{GridFitFilled}\\
d. $\dom g=#4$, $\rng g=#5$%
}

\begin{document}
\maketitle
\emph{Note: In the first two sections, the given graph is shown in gray.}
\section{Single-Step Transformations}
In each figure below, the given graph is the graph of $y=f(x)$. Graph the function in the question on the same axes. You may find it helpful to use the marked points as reference points.
\begin{smenum}
\mitemxx
{$g(x)=2f(x)$\\
\begin{GridSquareFitFilled}(2in)
\begin{scope}[black!50!white]
	\draw[thick] (-2,-1) -- (2,2) -- (4,-2);
	\foreach \x/\y in {-2/-1,2/2,4/-2}
		\Cdot{\x,\y};
\end{scope}
	\draw[thick] (-2,-2) -- (2,4) -- (4,-4);
	\foreach \x/\y in {-2/-1,2/2,4/-2}
		\Cdot{\x,2*\y};
\end{GridSquareFitFilled}}
{$g(x)=f(-x)$\\
\begin{GridSquareFitFilled}(2in)
\begin{scope}[black!50!white]
	\draw[thick] (-4,3) -- (3,-4);
	\foreach \x/\y in {-4/3,3/-4}
		\Cdot{\x,\y};
\end{scope}
	\draw[thick] (4,3) -- (-3,-4);
	\foreach \x/\y in {-4/3,3/-4}
		\Cdot{-\x,\y};
\end{GridSquareFitFilled}}
\mitemxx
{$g(x)=f(x)-3$\\
\begin{GridSquareFitFilled}(2in)
\begin{scope}[black!50!white]
	\draw[thick] (-3,0) arc (180:0:3);
	\foreach \x/\y in {-3/0,0/3,3/0}
		\Cdot{\x,\y};
\end{scope}
	\draw[thick] (-3,-3) arc (180:0:3);
	\foreach \x/\y in {-3/0,0/3,3/0}
		\Cdot{\x,\y-3};
\end{GridSquareFitFilled}}
{$g(x)=-f(x)$\\
\begin{GridSquareFitFilled}(2in)
\begin{scope}[black!50!white]
	\draw[thick] (4,4) -- (0,0) -- (-3,3);
	\foreach \x/\y in {-3/3,0/0,4/4}
		\Cdot{\x,\y};
\end{scope}
	\draw[thick] (4,-4) -- (0,0) -- (-3,-3);
	\foreach \x/\y in {-3/3,0/0,4/4}
		\Cdot{\x,-\y};
\end{GridSquareFitFilled}}
\mitemxx
{$g(x)=f(x)+2$\\
\begin{GridSquareFitFilled}(2in)
\begin{scope}[black!50!white]
	\draw[thick] (-3,3) parabola bend (1,-1) (3,0);
	\foreach \x/\y in {-3/3,1/-1,3/0}
		\Cdot{\x,\y};
\end{scope}
	\draw[thick] (-3,5) parabola bend (1,1) (3,2);
	\foreach \x/\y in {-3/3,1/-1,3/0}
		\Cdot{\x,\y+2};
\end{GridSquareFitFilled}}
{$g(x)=f(x+1)$\\
\begin{GridSquareFitFilled}(2in)
\begin{scope}[black!50!white]
	\draw[thick] (-3,-3) -- (-1,1) -- (3,1) -- (4,2);
	\foreach \x/\y in {-3/-3,-1/1,3/1,4/2}
		\Cdot{\x,\y};
\end{scope}
	\draw[thick] (-4,-3) -- (-2,1) -- (2,1) -- (3,2);
	\foreach \x/\y in {-3/-3,-1/1,3/1,4/2}
		\Cdot{\x-1,\y};
\end{GridSquareFitFilled}}
\mitemxx
{$g(x)=\frac{2}{3}f(x)$\\
\begin{GridSquareFitFilled}(2in)
\begin{scope}[black!50!white]
	\draw[thick] (-1,-3) -- (4,3);
	\foreach \x/\y in {-1/-3,4/3}
		\Cdot{\x,\y};
\end{scope}
	\draw[thick] (-1,-2) -- (4,2);
	\foreach \x/\y in {-1/-3,4/3}
		\Cdot{\x,2/3*\y};
\end{GridSquareFitFilled}}
{$g(x)=f(\frac{1}{3}x)$\\
\begin{GridSquareFitFilled}(2in)
\begin{scope}[black!50!white]
	\draw[thick] (-1,-2) parabola bend (0,3) (1,-2);
	\foreach \x/\y in {-1/-2,0/3,1/-2}
		\Cdot{\x,\y};
\end{scope}
	\draw[thick] (-3,-2) parabola bend (0,3) (3,-2);
	\foreach \x/\y in {-1/-2,0/3,1/-2}
		\Cdot{3*\x,\y};
\end{GridSquareFitFilled}}
\mitemxx
{$g(x)=f(2x)$\\
\begin{GridSquareFitFilled}(2in)
\begin{scope}[black!50!white]
	\draw[thick] (-4,1) arc (-180:0:4);
	\foreach \x/\y in {-4/1,0/-3,4/1}
		\Cdot{\x,\y};
\end{scope}
	\draw[thick,xscale=0.5] (-4,1) arc (-180:0:4);
	\foreach \x/\y in {-4/1,0/-3,4/1}
		\Cdot{\x/2,\y};
\end{GridSquareFitFilled}}
{$g(x)=f(x-3)$\\
\begin{GridSquareFitFilled}(2in)
\begin{scope}[black!50!white]
	\draw[thick] (-3,-2) -- (1,4);
	\foreach \x/\y in {-3/-2,1/4}
		\Cdot{\x,\y};
\end{scope}
	\draw[thick] (0,-2) -- (4,4);
	\foreach \x/\y in {-3/-2,1/4}
		\Cdot{\x+3,\y};
\end{GridSquareFitFilled}}
\end{smenum}

\section{Multiple Transformations}
In each figure below, the given graph is the graph of $y=f(x)$. Graph the function in the question on the same axes. You may find it helpful to use the marked points as reference points. \emph{You will want to write out your work on a separate sheet of paper.}
\begin{smenum}
\mitemxx
{$g(x)=-2f(\frac{2}{3}x+2)+3$\\
\begin{GridSquareFitFilled}[11]
\begin{scope}[black!50!white]
	\draw[thick] (-4,-4) -- (2,3) -- (6,-2);
	\Cdot{-4,-4};
	\Cdot{2,3};
	\Cdot{6,-2};
\end{scope}
	\draw[thick] (-9,11) -- (0,-3) -- (6,7);
	\Cdot{-9,11};
	\Cdot{0,-3};
	\Cdot{6,7};
\end{GridSquareFitFilled}}
%\mitemx
{$g(x)=\frac{2}{3}f(-2x-5)-1$\\
\begin{GridSquareFitFilled}[11]
\begin{scope}[black!50!white]
	\draw[thick] (-9,-3) arc (180:0:6);
	\Cdot{-9,-3};
	\Cdot{-3,3};
	\Cdot{3,-3};
\end{scope}
	\draw[thick,xscale=0.75] (-5.25,-3) arc (180:0:4);
	\Cdot{2,-3};
	\Cdot{-1,1};
	\Cdot{-4,-3};
\end{GridSquareFitFilled}}
\end{smenum}
\clearpage
\section{Domain and Range}
\begin{senum}
\item {$\dom g=\rint{1}$}
\item {$\dom h=\Rint{-6}\cup\Lint{-\frac{16}{3}}$}
\item {$\dom k = \lrint{15}{\frac{35}{2}}$, $\rng k=\eint$}
\end{senum}
\section{Symmetry}
Classify each function as even, odd, or neither.
\begin{smenum}
\mitemxxxx
{Even}
{Neither}
{Odd}
{Odd}
\mitemxxxx
{Even}
{Neither}
{Neither}
{Odd}
\mitemxxxx
{Neither}
{Odd}
{Even}
{Neither}
\mitemx
{A polynomial is an even function if all powers of $x$ are even (including a possible constant). It is odd if all powers of $x$ are odd (and there is no constant). If there is a mix, it is neither.}
\end{smenum}
\section{Using the Library}
For each function below:
\begin{clist}
\item Identify the base function (call it $f(x)$) and the family of functions it belongs to.
\item Rewrite the function in terms of $f(x)$. (We usually won't do this explicitly in practice.)
\item Sketch a graph of the function, including any asymptotes and clearly marking any special points, using a separate sheet of graph paper or the other side of this page.
\item Give the domain and range of the function.
\end{clist}
\begin{smenum}
\mitemxx
%%%%%%%%%%%%%%%%%%%%%%%%%%%%%%%%%%%%%%%%%%%%%%%%%%%%%%%%%%%%%%%%%%%%%%%%%%%%%%%%%%%%%%%%%%%%%%%
{\begin{fpa}{x}{Linear}{-3f(x)+5}{\eint}{\eint}(5,5,2,8)
\lineslopeintercept(-5,-2)(5,8){5}{-3}
\end{fpa}}
%%%%%%%%%%%%%%%%%%%%%%%%%%%%%%%%%%%%%%%%%%%%%%%%%%%%%%%%%%%%%%%%%%%%%%%%%%%%%%%%%%%%%%%%%%%%%%%
{\begin{fpa}{\dfrac{1}{x^2}}{Power (negative even)}{\dfrac{5}{(x-3)^2}=5f(x-3)}{\neint{3}}{\lint{0}}(5,5,2,8)
\draw[<->,thick,domain=-5.9:{3-sqrt(5/8.9)}] plot (\x,{5/(\x-3)^2});
\draw[<->,thick,domain={3+sqrt(5/8.9)}:5.9] plot (\x,{5/(\x-3)^2});
\draw[dashed,thick,yshift=1pt] (-5.9,0) -- (5.9,0);
\draw[dashed,thick] (3,-2.9) -- (3,8.9);
\end{fpa}}
\mitemxx
%%%%%%%%%%%%%%%%%%%%%%%%%%%%%%%%%%%%%%%%%%%%%%%%%%%%%%%%%%%%%%%%%%%%%%%%%%%%%%%%%%%%%%%%%%%%%%%
{\begin{fpa}{\dfrac{1}{x}}{Power (negative odd)}{3f(x-2)}{\neint{2}}{\neint{0}}%
\draw[<->,thick,domain=-5.9:{2-3/5.9}] plot (\x,{3/(\x-2)});
\draw[<->,thick,domain={2+3/5.9}:5.9] plot (\x,{3/(\x-2)});
\draw[dashed,thick,yshift=1pt] (-5.9,0) -- (5.9,0);
\draw[dashed,thick] (2,-5.9) -- (2,5.9);
\end{fpa}}
%%%%%%%%%%%%%%%%%%%%%%%%%%%%%%%%%%%%%%%%%%%%%%%%%%%%%%%%%%%%%%%%%%%%%%%%%%%%%%%%%%%%%%%%%%%%%%%
{\begin{fpa}{\Abs{x}}{Absolute value}{f(x-5)+2}{\eint}{\Lint{2}}(3,7,3,7)%
\draw[<->,thick] (-0.9,7.9) -- (5,2) -- (7.9,4.9);
\end{fpa}}
\mitemxx
%%%%%%%%%%%%%%%%%%%%%%%%%%%%%%%%%%%%%%%%%%%%%%%%%%%%%%%%%%%%%%%%%%%%%%%%%%%%%%%%%%%%%%%%%%%%%%%
{\begin{fpa}{1}{Constant}{4f(x)}{\eint}{\Set{4}}%
\draw[<->,thick] (-5.9,4) -- (5.9,4);
\end{fpa}}
%%%%%%%%%%%%%%%%%%%%%%%%%%%%%%%%%%%%%%%%%%%%%%%%%%%%%%%%%%%%%%%%%%%%%%%%%%%%%%%%%%%%%%%%%%%%%%%
{\begin{fpa}{x^3}{Power (odd positive)}{2f(x-1)+1}{\eint}{\eint}(1,3,5,5)%
\draw[<->,thick,domain={1-3.45^(1/3)}:{1+2.45^(1/3)}] plot (\x,{2*(\x-1)^3+1});
\end{fpa}}
%%%%%%%%%%%%%%%%%%%%%%%%%%%%%%%%%%%%%%%%%%%%%%%%%%%%%%%%%%%%%%%%%%%%%%%%%%%%%%%%%%%%%%%%%%%%%%%
\mitemxx
%%%%%%%%%%%%%%%%%%%%%%%%%%%%%%%%%%%%%%%%%%%%%%%%%%%%%%%%%%%%%%%%%%%%%%%%%%%%%%%%%%%%%%%%%%%%%%%
{\begin{fpa}{\log_{1/3}(x)}{Logarithm}{f(2x+3)+1}{\lint{-\frac{3}{2}}}{\eint}(2,8,5,5)%
\draw[<->,thick,domain={-ln(20.8)/ln(3)+1}:5.9] plot ({pow(1/3,\x-1)/2-3/2},\x);
\draw[thick,dashed] (-1.5,-5.9) -- (-1.5,5.9);
\end{fpa}}
%%%%%%%%%%%%%%%%%%%%%%%%%%%%%%%%%%%%%%%%%%%%%%%%%%%%%%%%%%%%%%%%%%%%%%%%%%%%%%%%%%%%%%%%%%%%%%%
{\begin{fpa}{\left(\frac{1}{2}\right)^x}{Exponential}{-f(x+2)+1}{\eint}{\rint{1}}(5,5,8,2)%
\draw[<->,thick,domain={ln(9.9)/ln(1/2)-2}:5.9] plot (\x,{-pow(1/2,\x+2)+1});
\draw[thick,dashed] (-5.9,1) -- (5.9,1);
\end{fpa}}
%%%%%%%%%%%%%%%%%%%%%%%%%%%%%%%%%%%%%%%%%%%%%%%%%%%%%%%%%%%%%%%%%%%%%%%%%%%%%%%%%%%%%%%%%%%%%%%
\mitemxx
%%%%%%%%%%%%%%%%%%%%%%%%%%%%%%%%%%%%%%%%%%%%%%%%%%%%%%%%%%%%%%%%%%%%%%%%%%%%%%%%%%%%%%%%%%%%%%%
{\begin{fpa}{x}{Linear}{2f(x-2)+1}{\eint}{\eint}%
\linepointslope(-5,-5)(5,5)(2,1){2}
\end{fpa}}
%%%%%%%%%%%%%%%%%%%%%%%%%%%%%%%%%%%%%%%%%%%%%%%%%%%%%%%%%%%%%%%%%%%%%%%%%%%%%%%%%%%%%%%%%%%%%%%
{\begin{fpa}{x^2}{Power (even positive)}{-3f(x+4)+2}{\eint}{\Rint{2}}(6,2,7,3)%
\draw[<->,thick] ({-4-sqrt(3.3)},-7.9) parabola bend(-4,2) ({-4+sqrt(3.3)},-7.9);
\end{fpa}}
%%%%%%%%%%%%%%%%%%%%%%%%%%%%%%%%%%%%%%%%%%%%%%%%%%%%%%%%%%%%%%%%%%%%%%%%%%%%%%%%%%%%%%%%%%%%%%%
\mitemxx
%%%%%%%%%%%%%%%%%%%%%%%%%%%%%%%%%%%%%%%%%%%%%%%%%%%%%%%%%%%%%%%%%%%%%%%%%%%%%%%%%%%%%%%%%%%%%%%
{\begin{fpa}{\dfrac{1}{x^2}}{Power (even negative)}{\frac{3}{2}f(x)-1}{\neint{0}}{\lint{-1}}%
\draw[<->,thick,domain=-5.9:{-pow(4.6,-1/2)}] plot (\x,{3/2*pow(\x,-2)-1});
\draw[<->,thick,domain={pow(4.6,-1/2)}:5.9] plot (\x,{3/2*pow(\x,-2)-1});
\draw[dashed,thick] (-5.9,-1) -- (5.9,-1);
\draw[dashed,thick,xshift=-1pt] (0,-5.9) -- (0,5.9);
\end{fpa}}
%%%%%%%%%%%%%%%%%%%%%%%%%%%%%%%%%%%%%%%%%%%%%%%%%%%%%%%%%%%%%%%%%%%%%%%%%%%%%%%%%%%%%%%%%%%%%%%
{\begin{fpa}{\sqrt[3]{x}}{Power (odd root)}{3f(4-x)-1}{\eint}{\eint}(4,6,5,5)%
\draw[<->,thick,domain={-3*pow(2.9,1/3)-1}:{3*pow(8.9,1/3)-1}] plot ({4-pow(\x/3+1/3,3)},\x);
\end{fpa}}
%%%%%%%%%%%%%%%%%%%%%%%%%%%%%%%%%%%%%%%%%%%%%%%%%%%%%%%%%%%%%%%%%%%%%%%%%%%%%%%%%%%%%%%%%%%%%%%
\mitemxx
%%%%%%%%%%%%%%%%%%%%%%%%%%%%%%%%%%%%%%%%%%%%%%%%%%%%%%%%%%%%%%%%%%%%%%%%%%%%%%%%%%%%%%%%%%%%%%%
{\begin{fpa}{\dfrac{1}{x}}{Power (odd negative)}{\dfrac{2}{2x-3}+3=2f(2x-3)+3}{\neint{\frac{3}{2}}}{\neint{3}}(3,7,3,7)%
\draw[<->,thick,domain=-3.9:18.4/13.8] plot (\x,{(6*\x-7)/(2*\x-3)});
\draw[<->,thick,domain=16.7/9.8:7.9] plot (\x,{(6*\x-7)/(2*\x-3)});
\draw[dashed,thick] (-3.9,3) -- (7.9,3);
\draw[dashed,thick,xshift=-1pt] (1.5,-3.9) -- (1.5,7.9);
\end{fpa}}
%%%%%%%%%%%%%%%%%%%%%%%%%%%%%%%%%%%%%%%%%%%%%%%%%%%%%%%%%%%%%%%%%%%%%%%%%%%%%%%%%%%%%%%%%%%%%%%
{\begin{fpa}{\sqrt{x}}{Power (even root)}{2f(-\frac{1}{2}x+3)}{\Rint{6}}{\Lint{0}}(3,7,3,7)%
\draw[->,thick,domain=0:{2*sqrt(4.95)}] plot (6-\x^2/2,\x);
\Cdot{6,0};
\end{fpa}}
%%%%%%%%%%%%%%%%%%%%%%%%%%%%%%%%%%%%%%%%%%%%%%%%%%%%%%%%%%%%%%%%%%%%%%%%%%%%%%%%%%%%%%%%%%%%%%%
\mitemxx
%%%%%%%%%%%%%%%%%%%%%%%%%%%%%%%%%%%%%%%%%%%%%%%%%%%%%%%%%%%%%%%%%%%%%%%%%%%%%%%%%%%%%%%%%%%%%%%
{\begin{fpa}{x^2}{Power (even positive)}{2(x-2)^2-3=2f(x-2)-3}{\eint}{\Lint{-3}}(1,5,4,6)%
\draw[<->,thick] ({2-sqrt(4.95)},6.9) parabola bend(2,-3) ({2+sqrt(4.95)},6.9);
\end{fpa}}
%%%%%%%%%%%%%%%%%%%%%%%%%%%%%%%%%%%%%%%%%%%%%%%%%%%%%%%%%%%%%%%%%%%%%%%%%%%%%%%%%%%%%%%%%%%%%%%
{\begin{fpa}{3^x}{Exponential}{f(\frac{1}{3}x)-5}{\eint}{\Lint{-5}}%
\draw[<->,thick,domain=-5.9:5.9] plot (\x,{pow(3,\x/3)-5});
\draw[dashed,thick] (-5.9,-5) -- (5.9,-5);
\Cdot{6,0};
\end{fpa}}
%%%%%%%%%%%%%%%%%%%%%%%%%%%%%%%%%%%%%%%%%%%%%%%%%%%%%%%%%%%%%%%%%%%%%%%%%%%%%%%%%%%%%%%%%%%%%%%
\mitemxx
%%%%%%%%%%%%%%%%%%%%%%%%%%%%%%%%%%%%%%%%%%%%%%%%%%%%%%%%%%%%%%%%%%%%%%%%%%%%%%%%%%%%%%%%%%%%%%%
{\begin{fpa}{\log_2(x)}{Logarithm}{f(x-3)+4}{\lint{3}}{\eint}(1,9,5,5)%
\draw[<->,thick,domain=-5.9:5.9] plot ({pow(2,\x-4)+3},\x);
\draw[dashed,thick](3,-5.9) -- (3,5.9);
\end{fpa}}
%%%%%%%%%%%%%%%%%%%%%%%%%%%%%%%%%%%%%%%%%%%%%%%%%%%%%%%%%%%%%%%%%%%%%%%%%%%%%%%%%%%%%%%%%%%%%%%
{\begin{fpa}{\Abs{x}}{Absolute value}{-f(4-\frac{1}{3}x)+2}{\eint}{\Rint{2}}(3,23,4,4)[2]%
\draw[<->,thick] (-3.9,-33/10) -- (12,2) -- (23.9,-59/30);
\end{fpa}}
%%%%%%%%%%%%%%%%%%%%%%%%%%%%%%%%%%%%%%%%%%%%%%%%%%%%%%%%%%%%%%%%%%%%%%%%%%%%%%%%%%%%%%%%%%%%%%%
\end{smenum}
\end{document}